\chapter{Manual Técnico}
\label{sec:manual-tecnico}

Este manual está dirigido a quien tenga que levantar SIGA en un entorno de desarrollo, mantenerlo o extenderlo. Es una guía operativa: cómo poner el sistema a correr en una máquina nueva, cómo agregar funcionalidad siguiendo las convenciones que ya usa el proyecto, y qué particularidades del código conviene conocer antes de tocarlo.

Deliberadamente no repite lo que ya está en el cuerpo del documento. Para entender \emph{por qué} el sistema está construido como está, la referencia es:

\begin{itemize}
  \item \textbf{Arquitectura} (capas, pipeline de seguridad, aislamiento multi-sucursal, servicios de fondo): Capítulo~\ref{sec:arch}.
  \item \textbf{Arquitectura del frontend} (estado, enrutado, permisos en la UI): Capítulo~\ref{sec:frontend}.
  \item \textbf{Modelo de datos} (entidades, relaciones, enums, convenciones de tipos): Capítulo~\ref{sec:modelo-datos}.
  \item \textbf{Despliegue en producción} (servidor, \code{docker-compose.}\allowbreak \code{prod.}\allowbreak \code{yml}, variables de entorno, TLS): Capítulo~\ref{sec:implementacion}.
  \item \textbf{Endpoints} (los 234, con su policy y sus DTOs): Apéndice~\ref{sec:api-reference}.
  \item \textbf{Decisiones de diseño} y su justificación: Apéndice~\ref{sec:adr}.
\end{itemize}

\section{Entorno de desarrollo}
\label{sec:mt-entorno}

\subsection{Requisitos}
\label{sec:mt-requisitos}

\begin{itemize}
  \item SDK de .NET 10, para compilar y correr el backend y para crear migraciones.
  \item Node \code{^20.19.0 || >=22.12.0}, para el frontend.
  \item Docker, para la base de datos local.
\end{itemize}

A diferencia de producción --- donde todo corre en contenedores y el host solo necesita Docker (ver Capítulo~\ref{sec:implementacion}) --- en desarrollo conviene tener el SDK y Node instalados en la máquina, para aprovechar la recarga en caliente de ambos lados.

\subsection{Levantar el sistema}
\label{sec:mt-levantar}

Los dos repositorios (\code{SIGA} y \code{SIGA-Web}) se clonan como directorios hermanos. Después, tres pasos, cada uno en su terminal:

\newpage
\begin{lstlisting}[numbers=none]
# 1. Base de datos: Postgres 16 en un contenedor
cd SIGA && docker compose up -d

# 2. Backend: escucha en http://localhost:5038
dotnet run --project SIGA/src/SIGA.Api

# 3. Frontend: Vite en modo desarrollo
cd SIGA-Web && npm install && npm run dev
\end{lstlisting}

El contenedor de la base se llama \code{siga-postgres} y publica el puerto \textbf{5433} del host (mapeo \code{5433:5432}), no el 5432 --- así no choca con un Postgres ya instalado localmente.

Al arrancar, el backend aplica las migraciones pendientes y siembra los datos base (permisos, roles y la cuenta administradora inicial) sin intervención manual:

\begin{lstlisting}[numbers=none]
await db.Database.MigrateAsync();
await DbSeeder.SeedAsync(db);
\end{lstlisting}

El frontend no necesita configuración de CORS: el proxy de Vite (\code{vite.config.ts}) reenvía \code{/api/*} y \code{/uploads/*} al backend (\code{VITE_API_URL}, por defecto \code{http:/}\allowbreak \code{/}\allowbreak \code{localhost:5038}). Por eso todos los servicios del frontend usan rutas relativas, y ese mismo esquema se reproduce en producción con el proxy de borde.

Scripts disponibles en el frontend:

\begin{lstlisting}[numbers=none]
npm run dev      # levanta Vite en modo desarrollo
npm run build    # type-check (vue-tsc --build) EN PARALELO con vite build
npm run preview  # sirve el dist ya construido
npm run format   # prettier --write sobre src/
\end{lstlisting}

\section{Dónde está cada cosa}
\label{sec:mt-estructura}

El mapa del código ya está en el cuerpo del documento y no se repite aquí, para que exista en un solo lugar: los cuatro proyectos del backend (\code{SIGA.}\allowbreak \code{Domain}, \code{SIGA.}\allowbreak \code{Application}, \code{SIGA.}\allowbreak \code{Infrastructure}, \code{SIGA.Api}), con qué contiene cada uno y en qué orden dependen entre sí, están en el Capítulo~\ref{sec:arch}, \S\ref{sec:arch-capas}; el árbol de \code{SIGA-Web/src/} (vistas, componentes, servicios, composables, stores, router, configuración del menú) está en el Capítulo~\ref{sec:frontend}, \S\ref{sec:frontend-estructura}.

Vale la pena tener presente una sola cosa al buscar dónde hacer un cambio: la lógica de negocio vive íntegramente en \code{SIGA.}\allowbreak \code{Infrastructure/}\allowbreak \code{Services/}. Los controllers no la tienen, y las entidades de \code{SIGA.Domain} solo tienen las reglas que dependen de sí mismas.

\section{Trabajar con migraciones}
\label{sec:mt-migraciones}

Crear una migración nueva:

\begin{lstlisting}[numbers=none]
dotnet ef migrations add <Nombre> --project SIGA/src/SIGA.Infrastructure --startup-project SIGA/src/SIGA.Api
\end{lstlisting}

Se aplican solas en el siguiente arranque de la API. Tres cosas a tener en cuenta:

\begin{itemize}
  \item \textbf{El prefijo \code{NNN} del nombre no indica el orden.} Hay 90 migraciones y el \code{NNN} más alto es 066: la numeración no es una secuencia completa, hay prefijos duplicados (003, 008, 021 a 027, 031 a 033, 041 a 043, 047, 048 y 059 aparecen dos veces cada uno) y 7 migraciones tempranas no llevan prefijo. El único orden real es el \emph{timestamp} de 14 dígitos con el que empieza el nombre del archivo, que es también el que usa el historial de migraciones de EF.
  \item \textbf{Con más de una réplica, todas ejecutan \code{MigrateAsync} a la vez} y compiten por aplicar la misma migración. Conviene arrancar de a una la primera vez, o correr la migración como paso separado del despliegue.
  \item \textbf{Conviene leer la migración autogenerada antes de aplicarla} y probarla en limpio, subida y bajada: a veces arrastra cambios no relacionados que quedaron pendientes en el modelo.
\end{itemize}

\section{Cómo agregar un endpoint nuevo}
\label{sec:mt-endpoint-nuevo}

Receta paso a paso, con los archivos a tocar en orden:

\begin{enumerate}
  \item \textbf{DTOs} en \code{SIGA.}\allowbreak \code{Application/}\allowbreak \code{DTOs/}\allowbreak \code{<Area>/}: el request (\code{Crear{Entity}Request}) y la response (\code{{Entity}Response} o \code{{Entity}Dto}). Nunca exponer la entidad de Domain.
  \item \textbf{Método en la interfaz} del servicio, en \code{SIGA.}\allowbreak \code{Application/}\allowbreak \code{Interfaces/}\allowbreak \code{I{Entity}Service.cs}, devolviendo \code{Task<Result<T>>} (o \code{Task<Result<Paged}\allowbreak \code{Result<T>>>} para listados).
  \item \textbf{Implementación} en \code{SIGA.}\allowbreak \code{Infrastructure/}\allowbreak \code{Services/}\allowbreak \code{{Entity}Service.cs}. Ahí va la lógica: validar toda la entrada y devolver \code{Result.Failure(msg, ErrorType.X)} en cada caso de error. Si la operación persiste stock, dinero u otra entidad con dimensión de sucursal, fijar \code{SucursalId} con \code{SucursalResolver.}\allowbreak \code{WriteBranchAsync}. \textbf{Para leer no hace falta filtrar por sucursal a mano}: el filtro global de EF ya lo hace (ver Capítulo~\ref{sec:arch}, \S\ref{sec:arch-multisucursal}).
  \item \textbf{Registrar el servicio}, si es nuevo, en \code{SIGA.}\allowbreak \code{Infrastructure/}\allowbreak \code{DependencyInjection.}\allowbreak \code{cs} dentro de \code{AddInfrastructure} (\code{services.}\allowbreak \code{AddScoped<}\allowbreak \code{I{Entity}Service, }\allowbreak \code{{Entity}Service>()}). Nunca en \code{Program.cs}.
  \item \textbf{Método en el controller} de \code{SIGA.}\allowbreak \code{Api/}\allowbreak \code{Controllers/}: liviano, decorado con \code{[Authorize(}\allowbreak \code{Policy = "...")]}, que llama al servicio y termina en \code{return }\allowbreak \code{ToHttpResponse(result)}. Sin lógica de negocio ni acceso al \code{AppDbContext}.
\end{enumerate}

\textbf{Un permiso nuevo se declara en dos lugares.} Si el endpoint usa un permiso que todavía no existe, hay que agregarlo tanto al arreglo \code{permission}\allowbreak \code{Policies} de \code{Program.cs} --- que es el que registra la \code{AuthorizationPolicy} de ASP.NET --- como a \code{DbSeeder.}\allowbreak \code{AllPermissions} y al rol correspondiente, que es lo que lo siembra en la base. Son dos listas separadas y nada las sincroniza: sembrar el permiso sin registrar la policy hace que el endpoint responda 500 (\emph{AuthorizationPolicy not found}) en vez de autorizar, y registrarla sin sembrar el permiso hace que nadie lo tenga y responda 403. La asimetría que existe hoy es \code{ver_recepcion}: está en el seed pero no tiene policy declarada, así que no puede usarse en un \code{[Authorize(Policy=.}\allowbreak \code{.}\allowbreak \code{.}\allowbreak \code{)]} aunque exista en la base.

\section{Cómo agregar una vista nueva}
\label{sec:mt-vista-nueva}

\begin{enumerate}
\item \textbf{Servicio} (\code{src/}\allowbreak \code{services/}\allowbreak \code{<entidad>Service.}\allowbreak \code{ts}). Conviene seguir el estilo del archivo vecino: si el área ya usa \code{useHttp}, usarlo; si no, \code{import { http } from "@/api/http"} directo. Se declaran las interfaces de request/response y una función por operación. Muchos listados devuelven \code{PagedResult<T>} (\code{items}, \code{totalCount}, \code{totalActive}, \code{page}, \code{pageSize}, \code{totalPages}) y reciben \texttt{\{ page?, pageSize?, search?, status? \}} serializados con \code{URLSearchParams}.

\item \textbf{Vista} (\code{src/}\allowbreak \code{views/}\allowbreak \code{<Nombre>View.}\allowbreak \code{vue}), con \code{<script setup lang="ts">}. Layout autenticado estándar:

\begin{lstlisting}[numbers=none]
<div class="min-h-screen" style="background-color: var(--color-background)">
  <AppSidebar />
  <AppHeader />
  <main style="margin-left: 280px; padding-top: 64px">
    <div class="p-8"><!-- contenido edge-to-edge, sin max-w ni mx-auto --></div>
  </main>
</div>
\end{lstlisting}

Se reusan los componentes base: encabezado con \code{<h1 class="text-4xl font-extrabold tracking-tight">} y \code{BaseButton} primary \code{lg}; \emph{toolbar} con \code{FilterChips} (izquierda) más \code{SearchInput class="w-72"} (derecha); \code{BaseTable} para el listado; \code{BaseModal} para crear/editar (\code{size="lg"}) y confirmar (\code{size="sm"}); inputs especializados (\code{MontoInput}, \code{DateInput}, \code{SearchableSelect}) en vez de nativos. Estilos de campo vía \code{inputStyle()} de \code{useFieldStyles}, nunca hex ni \code{font-family} inline. Los estados de carga, vacío y error se manejan con \code{try}/\code{catch} mapeando \code{err.status}.

\item \textbf{Ruta} en \code{src/}\allowbreak \code{router/}\allowbreak \code{index.}\allowbreak \code{ts}, con importación dinámica:

\begin{lstlisting}[numbers=none]
{
  path: "/mi-entidad",
  name: "mi-entidad",
  component: () => import("@/views/MiEntidadView.vue"),
  meta: { requiresAuth: true, permission: "ver_mi_entidad", label: "Mi Entidad" },
}
\end{lstlisting}

\item \textbf{Ítem de menú} en \code{src/}\allowbreak \code{config/}\allowbreak \code{menuConfig.}\allowbreak \code{ts}: en la \code{zone} que corresponda, con su \code{route}, su \code{icon} (Material Symbols) y su \code{permission}. Sin esto la vista existe pero no aparece en el \emph{sidebar}.

\item \textbf{Permisos de acción}: envolver los botones sensibles con \code{v-if="auth.}\allowbreak \code{hasPermission('.}\allowbreak \code{.}\allowbreak \code{.}\allowbreak \code{')"}, o usar el campo \code{hidden} de \code{RowContextMenu}.
\end{enumerate}

Antes de integrar el cambio: \code{npm run build} (que incluye el \emph{type-check}) y probar la pantalla logueado, en modo claro y en modo oscuro.

\section{Convenciones a tener en cuenta}
\label{sec:mt-convenciones}

Particularidades del código que no son evidentes al leerlo por primera vez y que conviene conocer antes de extenderlo:

\begin{itemize}
  \item \textbf{El error del backend no tiene una forma única.} \code{BaseController.}\allowbreak \code{ToHttpResponse} devuelve el error como string JSON pelado; \code{Empleados}\allowbreak \code{Controller} --- el único que no hereda de \code{BaseController} --- define su propio \code{ToResponse} que devuelve \texttt{\{ message \}}; y hay \emph{returns} ad-hoc con \texttt{\{ message \}} en algunos puntos. El frontend lo absorbe en un solo lugar, el interceptor de \code{src/api/http.ts}, con la cascada \code{typeof data === "string" ? data : (data?.message ?? data?.error ?? "Error {status}")}. Es tolerante, pero un endpoint nuevo que devuelva el detalle bajo otra clave (\code{detail}, \code{title}) hará que el usuario vea solo el código de estado. Al agregar endpoints conviene apuntar a \texttt{\{ message \}}.
  \item \textbf{Conviven dos estrategias de paginación.} \emph{Server-side}: 8 servicios devuelven \code{PagedResult<T>} y el pie de página usa \code{totalCount}/\code{totalPages} del backend. \emph{Client-side}: cerca de 18 vistas hacen \emph{slicing} local con una constante \code{PAGE_SIZE} (típicamente 10) sobre la lista completa ya traída. Antes de armar el pie de página de una pantalla conviene mirar cuál usa su servicio, sin asumir.
  \item \textbf{No hay contrato de paginación uniforme en la API.} \code{sortBy}/\code{sortOrder} no existen en ningún controller; el filtro de activo/inactivo no tiene estándar (\code{isActive} en unos, \code{status} string en otros, \code{soloActivos} en Empleados); y el \code{pageSize} por defecto varía por controller (10, 20 o 50).
  \item \textbf{La búsqueda por texto no usa índice.} Casi toda usa \code{.}\allowbreak \code{ToLower().}\allowbreak \code{Contains(q)} en EF, que se traduce a \texttt{lower(col) LIKE '\%q\%'}: con el comodín al principio el índice btree no sirve y termina en \emph{sequential scan}. Solo dos servicios usan \code{EF.}\allowbreak \code{Functions.}\allowbreak \code{ILike}. Cuando el volumen lo justifique, la salida es la extensión \code{pg_trgm} con índices GIN.
  \item \textbf{\code{BaseController} vive al final de \code{AuthController.cs}}, sin archivo propio. Es fácil no encontrarlo.
  \item \textbf{La carga de archivos tiene un solo caso y su patrón.} \code{ProductosController.}\allowbreak \code{UploadImagen} recibe un \code{IFormFile}, valida que no esté vacío ni supere 5\,MB y delega en \code{ProductoService.}\allowbreak \code{UploadImagenAsync}, que valida la extensión (jpg, jpeg, png, webp), guarda bajo \code{wwwroot/}\allowbreak \code{uploads/}\allowbreak \code{productos/} y persiste la ruta relativa en \code{Producto.}\allowbreak \code{ImagenUrl}. Se sirve como estático vía \code{app.}\allowbreak \code{UseStaticFiles()}. Es el único punto del sistema que guarda archivos en disco: conviene seguir ese patrón si aparece otro.
  \item \textbf{No hay pruebas automatizadas ni linter.} No existe proyecto de test en la solución, así que \code{dotnet test} no corre nada, y del lado del frontend no hay eslint. Los únicos controles automáticos son \code{dotnet build} y el \emph{type-check} de \code{vue-tsc}. Toda verificación funcional es manual: levantar el sistema, loguearse y recorrer el flujo.
  \item \textbf{Persona no es una sola tabla.} Coexisten \code{Patient} (clínico, con login) y \code{Cliente} (comercial, sin login) como dominios separados, cada uno con su controller y su servicio, colgando ambos de \code{Person}. El documento de identidad es \code{Person.CI} (cédula), no \code{DNI}. El razonamiento está en el \hyperref[adr:0001]{ADR 0001}.
\end{itemize}

\section{Construcción de las imágenes}
\label{sec:mt-imagenes}

El despliegue completo --- servidor, \code{docker-compose.}\allowbreak \code{prod.}\allowbreak \code{yml}, variables de entorno y TLS --- está en el Capítulo~\ref{sec:implementacion}. Acá quedan solo los detalles de construcción de cada imagen.

\textbf{Backend.} El \code{Dockerfile} es multi-etapa: compila con \code{mcr.}\allowbreak \code{microsoft.}\allowbreak \code{com/}\allowbreak \code{dotnet/}\allowbreak \code{sdk:10.}\allowbreak \code{0}, publica \code{SIGA.Api} y corre sobre \code{aspnet:10.0}, exponiendo el puerto 8080. Para construirlo suelto:

\begin{lstlisting}[numbers=none]
docker build -t siga:local SIGA/
\end{lstlisting}

\textbf{Frontend.} También multi-etapa: Node construye y nginx sirve el estático.

\begin{lstlisting}[numbers=none]
FROM node:20-alpine AS build
WORKDIR /app
COPY package*.json ./
RUN npm ci
COPY . .
ARG VITE_HCAPTCHA_SITE_KEY=10000000-ffff-ffff-ffff-000000000001
ENV VITE_HCAPTCHA_SITE_KEY=$VITE_HCAPTCHA_SITE_KEY
RUN npm run build

FROM nginx:alpine
COPY --from=build /app/dist /usr/share/nginx/html
COPY nginx-spa.conf /etc/nginx/conf.d/default.conf
EXPOSE 80
\end{lstlisting}

La \emph{site-key} de hCaptcha se resuelve en tiempo de \emph{build}: Vite reemplaza \code{import.}\allowbreak \code{meta.}\allowbreak \code{env.}\allowbreak \code{VITE_}\allowbreak \code{HCAPTCHA_}\allowbreak \code{SITE_}\allowbreak \code{KEY} por el valor literal en las pantallas públicas. El valor por defecto es la \emph{test key}, siempre válida; para el registro público real hay que pasar la de producción como \emph{build-arg}. No es un secreto: es una clave pública de cliente.

La segunda etapa usa \code{nginx-spa.conf}, cuyo único trabajo es hacer \emph{fallback} a \code{index.html} para que el enrutado del lado del cliente sobreviva a una recarga directa sobre una ruta interna:

\begin{lstlisting}[numbers=none]
server {
    listen 80;
    root /usr/share/nginx/html;
    index index.html;
    location / { try_files $uri $uri/ /index.html; }
}
\end{lstlisting}
